\subsection{应用}

命题 17. --- 设 G 为局部紧群，A 为 G 的一个子集，它对一个 Haar 测度可测且非局部可忽略。则 \(A \cdot A^{-1}\) 是 e 的一个邻域。

设 \(\beta\) 为一个左 Haar 测度。存在紧子集 \(K\)（\(G\) 的），使得 \(B = A \cap K\) 可积，测度 \(>0\)（对 \(\beta\)）。我们对 \(f = g = \varphi_{B}\) 应用命题 15 的推论。函数 \(F = \varphi_{B} * \check{\varphi}_{B}\) 连续且 \(>0\)（在 \(e\) 处）。因此存在一个邻域 \(V\)（\(e\) 的），使得 \(F(x) > 0\)（对 \(x \in V\)）。现在，

\[
\mathrm{F} (x) = \int \varphi_ {\mathrm{B}} (s) \varphi_ {\mathrm{B}} (x ^ {- 1} s) d \beta (s) = \beta (\mathrm{B} \cap x \mathrm{B}).
\]

因此，对 \(x \in V\)，有 \(B \cap xB \neq \varnothing\)，从而 \(x \in B \cdot B^{-1}\)。于是 \(V \subset B \cdot B^{-1} \subset A \cdot A^{-1}\)。

推论 1. --- 设 H 为 G 的一个子群，它对一个 Haar 测度 \(\beta\) 可测。则 H 或者是开的，或者是局部 \(\beta\)-可忽略的。

因为 \(H = H \cdot H^{-1}\)，所以若 H 不是局部 \(\beta\)-可忽略的，则 H 包含 e 的一个邻域（命题 17），从而是开的（GT, III, §2, No.~1 命题 4 的推论）。

推论 2. --- 设 L 为 G 的一个对乘法稳定的子集，其余集对一个 Haar 测度 \(\beta\) 局部可忽略。则 \(L = G\)。

因为 \(L^{-1}\) 与 \(L \cap L^{-1}\) 的余集都局部 \(\beta\)-可忽略。而 \(L \cap L^{-1}\) 是一个子群，故是开的（推论 1），从而是闭的。因此 \(G - (L \cap L^{-1})\)——它是开的且局部 \(\beta\)-可忽略的——是空的。于是 \(G = L \cap L^{-1}\)。

命题 18. --- 设 G 为局部紧群，\(\Gamma\) 为一个集合，其上带有一个乘法 \((u,v)\mapsto uv\) 与一个 Hausdorff 拓扑，使得：

\begin{enumerate}
\def\labelenumi{\arabic{enumi})}
\item
  \(\Gamma\) 的拓扑在平移下不变；
\item
  乘法在 \(\Gamma \times \Gamma\) 的每一紧子集上的限制是连续的。
\end{enumerate}

设 \(f: \mathbf{G} \to \Gamma\) 为 \(\mathbf{G}\) 到 \(\Gamma\) 中的一个映射，使得 \(f(xy) = f(x)f(y)\)（对 \(x, y\) 属于 \(\mathbf{G}\)），且对一个 Haar 测度 \(\beta\)（\(\mathbf{G}\) 上的）可测。则 \(f\) 连续。

置 \(g(x) = f(x^{-1})\)（对 \(x \in \mathbf{G}\)）。由于 \(f\) 与 \(g\) 都是 \(\beta\)-可测的，存在一个非 \(\beta\)-可忽略的紧子集 \(\mathbf{K}\)（\(\mathbf{G}\) 的），使得 \(f\) 与 \(g\) 在 \(\mathbf{K}\) 上的限制连续。映射 \((x,y) \mapsto f(xy^{-1}) = f(x)g(y)\)（从 \(\mathbf{K} \times \mathbf{K}\) 到 \(\Gamma\) 中）连续，因为 \(\Gamma\) 的乘法在 \(f(\mathbf{K}) \times g(\mathbf{K})\) 上连续；而这个映射可写成 \(\varphi \circ \psi\)，其中 \(\psi\) 是映射 \((x,y) \mapsto xy^{-1}\)（从 \(\mathbf{K} \times \mathbf{K}\) 到 \(\mathbf{K} \cdot \mathbf{K}^{-1}\) 上），\(\varphi\) 是 \(f\) 在 \(\mathbf{K} \cdot \mathbf{K}^{-1}\) 上的限制。设 \(\mathbf{R}\) 为在 \(\mathbf{K} \times \mathbf{K}\) 上由 \(\psi\) 定义的等价关系。映射 \(\psi'\)（\((\mathbf{K} \times \mathbf{K}) / \mathbf{R}\) 到 \(\mathbf{K} \cdot \mathbf{K}^{-1}\) 上的、由 \(\psi\) 过渡到商而导出的）连续，因此 \((\mathbf{K} \times \mathbf{K}) / \mathbf{R}\) 是 Hausdorff 的且 \(\psi'\) 是一个同胚。由于 \(\varphi \circ \psi\) 连续，可见 \(f\) 在 \(\mathbf{K} \cdot \mathbf{K}^{-1}\) 上的限制连续。而 \(\mathbf{K} \cdot \mathbf{K}^{-1}\) 是 \(e\) 的一个邻域（命题 17），因此 \(f\) 在 \(e\) 处连续。对每一 \(x_0 \in \mathbf{G}\)，\(f(x_0x) = f(x_0)f(x)\)，故 \(f\) 在 \(x_0\) 处连续，因为 \(\Gamma\) 的拓扑在平移下不变。

推论 1. --- 设 G 为局部紧群，\(\beta\) 为 G 上的一个 Haar 测度，E 为一个 Hausdorff 桶型局部凸空间，\(U\) 为 G 在 E 中的一个线性表示，使得 \(U(s) \in \mathcal{L}(\mathrm{E};\mathrm{E})\) 对所有 \(s \in \mathrm{G}\) 成立，且是 \(\beta\)-可测的（当 \(\mathcal{L}(\mathrm{E};\mathrm{E})\) 赋有点态收敛拓扑时）。则 \(U\) 是一个连续线性表示。

设 \(\Gamma\) 为 E 的自同构群，赋有点态收敛拓扑。这一拓扑是 Hausdorff 的且在平移下不变。设 K 为 \(\Gamma\) 的一个紧子集。则 K 在赋有点态收敛拓扑的 \(\mathcal{L}(\mathrm{E};\mathrm{E})\) 中有界，故是等度连续的（TVS, III, §4, No.~2 定理 1）；因此映射 \((u,v)\mapsto v\circ u\)（从 \(\mathbf{K}\times \mathbf{K}\) 到 \(\mathcal{L}(\mathbf{E};\mathbf{E})\) 中）连续（同处, §5, No.~5 命题 9 的推论 1）。于是，对每一 \(x\in \mathbf{E}\)，G 到 E 中的映射 \(s\mapsto U(s)x\) 连续（命题 18）。由于 E 是桶型的，\(U\) 连续（§2, No.~1 命题 1）。

推论 2. --- 设 G 为局部紧群，\(\beta\) 为 G 上的一个 Haar 测度，E 为一个可分 Banach 空间，\(U\) 为 G 在 E 中的一个线性表示，使得 \(U(s) \in \mathcal{L}(\mathrm{E};\mathrm{E})\) 对所有 \(s \in \mathrm{G}\) 成立。设 \((a_m)\) 为 E 中的一个完全序列，\((a_n')\) 为 E 的对偶 E' 的单位球 B'（赋弱拓扑）中的一个稠密序列。假设 G 上的函数 \(s \mapsto \langle U(s)a_m, a_n' \rangle\) 都是 \(\beta\)-可测的。则 \(U\) 是一个连续线性表示。

我们先证明，对每一 \(z' \in E'\)，纯量函数

\[
s \mapsto \langle U (s) a _ {m}, z ^ {\prime} \rangle
\]

是 \(\beta\)-可测的；我们可以限制在 \(\|z'\| \leqslant 1\) 的情形，并且由于 \(B'\) 对弱拓扑是可度量的（TVS, III, §3, No.~4 命题 6 的推论 2），存在一个子列 \((a_{n_k}')\)（\((a_n')\) 的），它弱收敛于 \(z'\)；函数

\[
s \mapsto \langle U (s) a _ {m}, z ^ {\prime} \rangle
\]

从而是一列 \(\beta\)-可测函数的极限，由此得出我们的断言。由此推出，映射 \(s \mapsto U(s)a_{m}\)（G 到 E 中的）对每个 m 是 \(\beta\)-可测的（Ch. IV, §5, No.~5 命题 10）。另一方面，存在 E 的元素的一个序列 \((b_{m})\)——它们是 \(a_{i}\) 的线性组合——在 E 的单位球中稠密。对每一 \(s \in \mathbf{G}\)，\(\|U(s)\| = \sup_{m} \|U(s)b_{m}\|\)，因此 \(s \mapsto \|U(s)\|\) 可测。设 K 为 G 的一个紧子集，且设 \(\varepsilon > 0\)。存在 K 的一个紧子集 \(\mathrm{K}_{0}\)，使得 \(\beta(\mathrm{K} - \mathrm{K}_{0}) \leqslant \varepsilon\)，且函数在 \(\mathrm{K}_{0}\) 上的限制连续——这里函数指 \(s \mapsto U(s)a_{m}\) 与 \(s \mapsto \|U(s)\|\)。于是诸 \(U(s)\)（对 \(s \in \mathrm{K}_{0}\)）是等度连续的，且点态收敛拓扑在 \(U(\mathrm{K}_{0})\) 上诱导出在 \(a_{m}\) 组成的集合上的点态收敛拓扑（TVS, III, §3, No.~4 命题 5）。因此映射 \(s \mapsto U(s)\)（从 \(\mathrm{K}_{0}\) 到 \(\mathcal{L}_{s}(\mathrm{E};\mathrm{E})\) 中）连续。然后只需应用推论 1。

